% 00-siglas.tex — siglas registradas na lista sem aparecer por extenso aqui (issue 98).
% As que aparecem por extenso no texto usam \criarsigla* no ponto da 1ª ocorrência
% (MLC, NDP, ABNT, CEN, NPT, CBMPR, CCA, CCB, BDI); EC5 é apresentada junto do CEN.
\criarsigla{EC5}{Eurocódigo 5, EN 1995-1-1: projeto de estruturas de madeira}
\criarsigla{EN}{norma europeia}
\criarsigla{NBR}{norma brasileira}
\criarsigla{ISO}{Organização Internacional de Normalização}
\criarsigla{UC}{classe de uso, EN 335}
\criarsigla{NP}{classe de penetração do preservativo, EN 351-1}
\criarsigla{CV-H}{contraventamento horizontal, da cobertura}
\criarsigla{CV-V}{contraventamento vertical, painéis das fachadas e do alinhamento central}
\criarsigla{CV-T}{contraventamento transversal, painéis dos frontões}
\criarsigla{CV-D}{alma do diafragma da cobertura, cruzes nas baías intermediárias}
\criarsigla{L1, L1c, L1f}{ligações de topo: viga sobre pilar lateral, sela do pilar central e topo do pilar de frontão}
\criarsigla{L3, L3c}{ligações de base correntes: pilar lateral e pilar central}
\criarsigla{L3f, L3a}{ligações de base do pilar de frontão e base ancorada dos montantes}
\criarsigla{P, PF, T, V}{pilares, pilares de frontão, terças e vigas, numerados conforme a tabela Nomenclatura dos elementos, que define também as designações acima}
